% GENERATED FILE. Edit canonical lessons, then run npm run generate:books.
% canonical-source-hash: fnv1a64:a20900b6d72ddee5
% canonical-lessons: RU-C233-reshat, RU-C233-vybirat, RU-C233-menyat, RU-C233-pokazyvat, RU-C233-prinosit

\chapter{To decide, To choose, To change, To show, To bring}
\label{ch:ru-to-decide-to-choose-to-change-to-show-to-bring}
\hlchaptermodality{233}

\begin{chapteropening}
\textbf{By the end of this chapter:} \emph{I can say to decide, to choose, to change, to show and to bring.}

Complete the last lesson of chapter 233: I can say to decide, to choose, to change, to show and to bring.
\end{chapteropening}

\section[\texorpdfstring{\ru{решать} (reshát)}{решать (reshát)}]{\ru{решать} (reshát) — to decide; to solve}

\label{lesson:RU-C233-reshat}



\begin{quote}
Before the new one: say the Russian for to send, then the Russian for to receive.
\end{quote}

\subsection*{You'll want to know: \ru{решать}}
\textbf{\ru{решать}} (\emph{reshát}) — \textquotedblleft{}to decide; to solve\textquotedblright{}.

\begin{grammarlens}[title={What you've built}]
One more word for this chapter.
\end{grammarlens}

\subsection*{Your turn}
\begin{itemize}
\raggedright
  \item \emph{Say it:} \emph{reshát}
  \item \emph{Say it:} \emph{reshát}, once more
  \item \emph{Say it:} say \emph{reshát}
  \item \emph{Recall:} say the Russian for to send, then the Russian for to receive, then say \emph{reshát} again
\end{itemize}

\begin{tcolorbox}[breakable,colback=teal!4,colframe=teal!35!black,title={Before you move on}]
What does \ru{решать} mean? (\textquotedblleft{}To decide; to solve\textquotedblright{}.) Say it once more.
\end{tcolorbox}
\section[\texorpdfstring{\ru{выбирать} (vybirát)}{выбирать (vybirát)}]{\ru{выбирать} (vybirát) — to choose}

\label{lesson:RU-C233-vybirat}



\begin{quote}
Before the new one: say the Russian for to receive, then the Russian for to decide; to solve.
\end{quote}

\subsection*{You'll want to know: \ru{выбирать}}
\textbf{\ru{выбирать}} (\emph{vybirát}) — \textquotedblleft{}to choose\textquotedblright{}.

\begin{grammarlens}[title={What you've built}]
One more word for this chapter.
\end{grammarlens}

\subsection*{Your turn}
\begin{itemize}
\raggedright
  \item \emph{Say it:} \emph{vybirát}
  \item \emph{Say it:} \emph{vybirát}, once more
  \item \emph{Say it:} say \emph{vybirát}
  \item \emph{Recall:} say the Russian for to receive, then the Russian for to decide; to solve, then say \emph{vybirát} again
\end{itemize}

\begin{tcolorbox}[breakable,colback=teal!4,colframe=teal!35!black,title={Before you move on}]
What does \ru{выбирать} mean? (\textquotedblleft{}To choose\textquotedblright{}.) Say it once more.
\end{tcolorbox}
\section[\texorpdfstring{\ru{менять} (menyát)}{менять (menyát)}]{\ru{менять} (menyát) — to change}

\label{lesson:RU-C233-menyat}



\begin{quote}
Before the new one: say the Russian for to decide; to solve, then the Russian for to choose.
\end{quote}

\subsection*{You'll want to know: \ru{менять}}
\textbf{\ru{менять}} (\emph{menyát}) — \textquotedblleft{}to change\textquotedblright{}.

\begin{grammarlens}[title={What you've built}]
One more word for this chapter.
\end{grammarlens}

\subsection*{Your turn}
\begin{itemize}
\raggedright
  \item \emph{Say it:} \emph{menyát}
  \item \emph{Say it:} \emph{menyát}, once more
  \item \emph{Say it:} say \emph{menyát}
  \item \emph{Recall:} say the Russian for to decide; to solve, then the Russian for to choose, then say \emph{menyát} again
\end{itemize}

\begin{tcolorbox}[breakable,colback=teal!4,colframe=teal!35!black,title={Before you move on}]
What does \ru{менять} mean? (\textquotedblleft{}To change\textquotedblright{}.) Say it once more.
\end{tcolorbox}
\section[\texorpdfstring{\ru{показывать} (pokázyvat)}{показывать (pokázyvat)}]{\ru{показывать} (pokázyvat) — to show}

\label{lesson:RU-C233-pokazyvat}



\begin{quote}
Before the new one: say the Russian for to choose, then the Russian for to change.
\end{quote}

\subsection*{You'll want to know: \ru{показывать}}
\textbf{\ru{показывать}} (\emph{pokázyvat}) — \textquotedblleft{}to show\textquotedblright{}.

\begin{grammarlens}[title={What you've built}]
One more word for this chapter.
\end{grammarlens}

\subsection*{Your turn}
\begin{itemize}
\raggedright
  \item \emph{Say it:} \emph{pokázyvat}
  \item \emph{Say it:} \emph{pokázyvat}, once more
  \item \emph{Say it:} say \emph{pokázyvat}
  \item \emph{Recall:} say the Russian for to choose, then the Russian for to change, then say \emph{pokázyvat} again
\end{itemize}

\begin{tcolorbox}[breakable,colback=teal!4,colframe=teal!35!black,title={Before you move on}]
What does \ru{показывать} mean? (\textquotedblleft{}To show\textquotedblright{}.) Say it once more.
\end{tcolorbox}
\section[\texorpdfstring{\ru{приносить} (prinosít)}{приносить (prinosít)}]{\ru{приносить} (prinosít) — to bring}

\label{lesson:RU-C233-prinosit}



\begin{quote}
Before the new one: say the Russian for to change, then the Russian for to show.
\end{quote}

\subsection*{You'll want to know: \ru{приносить}}
\textbf{\ru{приносить}} (\emph{prinosít}) — \textquotedblleft{}to bring\textquotedblright{}.

\begin{grammarlens}[title={What you've built}]
One more word for this chapter.
\end{grammarlens}

\subsection*{Your turn}
\begin{itemize}
\raggedright
  \item \emph{Say it:} \emph{prinosít}
  \item \emph{Say it:} \emph{prinosít}, once more
  \item \emph{Say it:} say \emph{prinosít}
  \item \emph{Recall:} say the Russian for to change, then the Russian for to show, then say \emph{prinosít} again
\end{itemize}

\begin{tcolorbox}[breakable,colback=teal!4,colframe=teal!35!black,title={Before you move on}]
What does \ru{приносить} mean? (\textquotedblleft{}To bring\textquotedblright{}.) Say it once more.
\end{tcolorbox}
